% SYNTHETIC EVALUATION DRAFT — written for the retrieval eval in evals/.
% Near-topic distractor for rl_trading.tex. Not a real paper; contains no results.
\documentclass{article}
\title{Population-Based Self-Play for Cooperative and Competitive Multi-Agent Reinforcement Learning in Games}
\begin{document}
\maketitle

\begin{abstract}
Multi-agent reinforcement learning (MARL) studies several agents that learn
simultaneously in a shared environment, cooperating, competing, or both. From
the perspective of any single agent the environment is non-stationary because
the other agents' policies keep changing. This draft proposes a
population-based self-play method with a centralized value function for mixed
cooperative-competitive games, and plans an evaluation on team-based
real-time strategy micromanagement, simulated football and imperfect
information card games.
\end{abstract}

\section{Introduction}
Many sequential decision problems involve several interacting decision makers:
teams of units in strategy games, players in sports simulations, and players
in poker or cooperative card games. MARL extends single-agent reinforcement
learning to a Markov game, in which the transition and reward depend on the
joint action of all agents. Solution concepts come from game theory: in
zero-sum games the target is a Nash equilibrium, while in cooperative games
the goal is to maximize a shared team reward.

Three challenges arise. Non-stationarity: each agent's learning changes the
environment for the others, which breaks the convergence assumptions of
independent learners. Credit assignment: in cooperative settings with a single
team reward, an agent cannot easily tell how its own action contributed.
Cycling and exploitability: naive self-play against the latest policy can
cycle between strategies in non-transitive games instead of improving.

Our contributions are: (1) a centralized-training, decentralized-execution
actor-critic in which each agent's policy uses only local observations, while a
centralized critic conditions on the global state and all actions; (2) a
population of policies trained with prioritized fictitious self-play, where
opponents are sampled from the population by their win rate against the
learner; and (3) an exploitability-based evaluation protocol.

\section{Related Work}
Independent Q-learning treats other agents as part of the environment.
Centralized training with decentralized execution underlies MADDPG, COMA, which
uses a counterfactual baseline for credit assignment, and value factorization
methods such as VDN and QMIX, which decompose the team value into per-agent
utilities. MAPPO showed that PPO with a centralized value function is a strong
baseline on cooperative benchmarks. In competitive games, fictitious play,
policy-space response oracles and league training with self-play reached
expert level in Go, poker and StarCraft. Emergent communication and opponent
modelling are related lines of work.

\section{Method}
\subsection{Centralized Critic}
Each agent $i$ has a policy $\pi_i(a_i \mid o_i)$ over its local observation.
The critic $V(s)$ takes the global state and is trained with generalized
advantage estimation; agents in the same team share the critic, and policies
are updated with the PPO clipped objective.

\subsection{Population and Opponent Sampling}
We maintain a population of past and current policies. For each training
episode, the opposing team is drawn from the population with probability
proportional to a function of the learner's loss rate against it, which
focuses training on opponents the learner cannot yet beat and prevents
forgetting.

\section{Experimental Plan}
We plan to evaluate on StarCraft multi-agent micromanagement scenarios, a
simulated football environment, and the cooperative card game Hanabi, plus a
two-team competitive gridworld. Baselines are independent PPO, MAPPO, QMIX and
naive self-play. Metrics are win rate, team return, and approximate
exploitability measured by training a best-response agent against frozen
policies.

\section{Limitations}
Population training multiplies compute. Exploitability can only be estimated
approximately in large games, and centralized critics assume the global state
is available during training.

\end{document}
